\chapter{Marco Teórico} \label{cap:marco_teorico}

Este capítulo presenta los conceptos sobre los que se apoya el \textit{pipeline}: los modelos de lenguaje pequeños y sus técnicas de compresión y adaptación, la tarea de detección de anomalías en logs, las métricas de evaluación, los \textit{embeddings} y la búsqueda vectorial, el uso de LLM como anotadores con medición de acuerdo, y la validación mediante inyección de fallos.

\section{Modelos de lenguaje pequeños}\label{sec:mt_slm}

\subsection{LLM y SLM}

Los modelos de lenguaje de gran tamaño (LLM), con decenas o cientos de miles de millones de parámetros, muestran capacidades de razonamiento y comprensión contextual, pero exigen recursos de cómputo y memoria elevados para su despliegue~\cite{wang2025_slm_survey, chen2025_collab_survey}. Los modelos de lenguaje pequeños (SLM) son modelos compactos, desde millones hasta unos pocos miles de millones de parámetros (por ejemplo, DistilBERT, TinyBERT, Llama-3.2-1B/3B o Phi-3.5), que buscan equilibrar capacidad y eficiencia en entornos de recursos limitados~\cite{wang2025_slm_survey, saha2025_slm_minimal}.
% TODO: definir con precisión el umbral de parámetros que usa esta tesis para llamar "pequeño" a un modelo (el estudio incluye modelos de 0.6B a 9B parámetros; Qwen3.5-9B es el más grande).

\subsection{Técnicas de compresión}\label{subsec:mt_compresion}

La revisión de Wang et al.~\cite{wang2025_slm_survey} agrupa las técnicas para obtener modelos pequeños y eficientes en tres familias:

\begin{description}
  \item[Cuantización.] Reduce la precisión numérica de pesos y activaciones (por ejemplo, de 16 bits a 8, 4 o menos bits), con ahorros de memoria de varios factores y pérdida moderada de calidad. Se distingue la cuantización posterior al entrenamiento de la cuantización consciente del entrenamiento.
  \item[Poda.] Elimina parámetros o componentes poco relevantes, ya sea pesos individuales (no estructurada) o estructuras completas como capas, canales o cabezas de atención (estructurada).
  \item[Destilación de conocimiento.] Transfiere el comportamiento de un modelo \textit{profesor} grande a un modelo \textit{estudiante} pequeño.
\end{description}

En este trabajo se emplean modelos ya cuantizados en formato GGUF (por ejemplo, Q4\_K\_M y Q6\_K) servidos con Ollama, junto con modelos en precisión completa cargados con Hugging Face Transformers (véase la Sección~\ref{sec:met_slm}).

\subsection{Ventajas del despliegue \textit{on-premise}}

Ejecutar el modelo dentro de la infraestructura institucional permite que la telemetría no salga de la organización, elimina el costo recurrente por \textit{token} de las API externas y ofrece latencias bajas adecuadas para análisis continuo~\cite{wang2025_slm_survey, pardeshi2026_smartllm, ma2026_adaptivelog}. % TODO: las cifras concretas de latencia y costo que reportan SmartLLM y AdaptiveLog están en notes/03_slm_fundamentos.md; incluirlas solo con la página/tabla exacta verificada.

\section{Ajuste fino eficiente en parámetros}\label{sec:mt_peft}

El ajuste fino completo de un modelo es costoso. Los métodos de ajuste fino eficiente en parámetros (PEFT) entrenan una fracción pequeña de parámetros adicionales manteniendo congelado el modelo base. En la literatura de análisis de logs se discuten tres:

\begin{itemize}
  \item \textbf{LoRA} (\textit{Low-Rank Adaptation}): congela los pesos preentrenados $W$ e inyecta matrices de bajo rango entrenables $A$ y $B$ tales que la actualización es $\Delta W = BA$, típicamente en las proyecciones de atención~\cite{lim2025_logadreft, wang2025_slm_survey}. % TODO: \cite{TODO} del artículo original de LoRA (Hu et al.), no está en el notebook ni en el .bib.
  \item \textbf{ReFT} (\textit{Representation Fine-Tuning}): en lugar de modificar pesos, aplica intervenciones de bajo rango sobre las representaciones ocultas entre capas; LogADReft lo aplica a detección de anomalías en logs~\cite{lim2025_logadreft}.
  \item \textbf{LINEADAPTER}: adaptador que inserta una capa convolucional 1D ligera tras la capa \textit{feed-forward} de los bloques \textit{Transformer} congelados, para capturar dependencias secuenciales locales en logs~\cite{liu2025_lineadapter}.
\end{itemize}

El \textit{fine-tuning} realizado en este trabajo utiliza LoRA (véase la Sección~\ref{sec:met_finetuning}).

\section{Detección de anomalías en logs}\label{sec:mt_anomalias}

\subsection{Formulación}

Un \textit{log} es un registro textual con marca de tiempo, nivel de severidad (\texttt{TRACE}, \texttt{DEBUG}, \texttt{INFO}, \texttt{WARN}, \texttt{ERROR}), origen (\textit{logger}) y mensaje, a veces acompañado de una traza de pila. La detección de anomalías consiste en decidir, para un evento o para un conjunto de eventos, si representa un comportamiento anómalo o una falla~\cite{he2021_survey}. En este trabajo la decisión se toma sobre un \textit{batch} de eventos consecutivos: el \textit{batch} es positivo (\emph{error}) si contiene al menos un evento que constituye una falla del sistema según el criterio operativo definido en la Sección~\ref{sec:met_criterio}.

\subsection{Enfoques}

Los enfoques se clasifican en supervisados (requieren etiquetas), no supervisados (modelan el comportamiento normal y marcan desviaciones) y basados en modelos de lenguaje, que pueden operar en modo \textit{zero-shot} o \textit{few-shot} mediante \textit{prompting}~\cite{liu2024_logprompt, karlsen2024_benchmark} o tras ajuste fino~\cite{lim2025_logadreft, liu2025_lineadapter}. Los costos de usar modelos grandes en todo el flujo han motivado esquemas colaborativos, donde un SLM atiende los casos sencillos y un LLM los difíciles~\cite{ma2026_adaptivelog}.

\subsection{Del diagnóstico a la causa raíz y la recomendación}

El análisis de causa raíz (RCA) busca identificar el componente o evento que origina una falla; las revisiones recientes cubren métodos basados en grafos, trazas y modelos de lenguaje~\cite{wang2024_rca}. Algunos trabajos combinan detección y RCA con LLM y redes bayesianas~\cite{pedroso2025_maia}. En el presente trabajo, la recomendación (causa probable y acción sugerida) se produce con un modelo de lenguaje local; no se evalúa formalmente su calidad frente a un conjunto de referencia.
% TODO: si se evalúa la calidad de las recomendaciones (p.ej. revisión experta de una muestra), agregarlo aquí y en Resultados.

\section{Métricas de evaluación}\label{sec:mt_metricas}

Sea $TP$, $FP$, $TN$ y $FN$ el número de verdaderos positivos, falsos positivos, verdaderos negativos y falsos negativos, tomando \emph{error} como clase positiva.

\begin{align}
\mathrm{Accuracy} &= \frac{TP + TN}{TP + FP + TN + FN}\\
\mathrm{Precision} &= \frac{TP}{TP + FP}\\
\mathrm{Recall} &= \frac{TP}{TP + FN}\\
\mathrm{Especificidad} &= \frac{TN}{TN + FP}\\
F_1 &= \frac{2\,\mathrm{Precision}\cdot\mathrm{Recall}}{\mathrm{Precision} + \mathrm{Recall}}
\end{align}

La curva ROC grafica la tasa de verdaderos positivos frente a la de falsos positivos al variar el umbral de decisión sobre una puntuación continua; el área bajo la curva (AUC) resume la capacidad de discriminación independientemente del umbral. En el \textit{pipeline}, la puntuación es la probabilidad $P(\text{error})$ que el modelo asigna al \textit{token} que escribe el valor booleano de su respuesta JSON, obtenida de los \textit{logits} (Hugging Face) o de los \textit{logprobs} (Ollama). El coeficiente de correlación de Matthews y la exactitud balanceada se reportan como complemento cuando la clase positiva es minoritaria.

Con clases desbalanceadas, \textit{accuracy} puede ser engañosa: un clasificador que nunca predice \emph{error} obtiene alta exactitud si los errores son raros. Por eso se reportan \textit{precision}, \textit{recall}, $F_1$ y AUC además de \textit{accuracy}.

\section{\textit{Embeddings} y búsqueda vectorial}\label{sec:mt_embeddings}

Un modelo de \textit{embeddings} transforma un texto en un vector denso en el que textos semánticamente similares quedan próximos (similitud coseno). Aplicado a logs, permite comparar mensajes por significado y no por coincidencia literal, y es más robusto ante vocabulario no visto y cambios de formato que los analizadores sintácticos~\cite{murphy2024_embeddings}. El proyecto utiliza el modelo \texttt{sentence-transformers/all-MiniLM-L6-v2} y una base vectorial Milvus Lite. Se emplean dos usos:

\begin{itemize}
  \item \textbf{Caché semántico:} si un \textit{batch} nuevo tiene similitud coseno mayor o igual a un umbral (0.8 por defecto) con uno ya clasificado, se reutiliza la clasificación.
  \item \textbf{Evaluación semántica:} cada \textit{batch} del SLM se compara con sus $k$ \textit{batches} más cercanos del \textit{ground truth} (vecinos más próximos, $k$-NN), lo que permite evaluar corridas con tamaños de \textit{batch} distintos al de la referencia.
\end{itemize}

\section{LLM como anotadores y acuerdo entre anotadores}\label{sec:mt_anotacion}

Cuando no existen etiquetas, una práctica es el \textit{pseudo-etiquetado}: varios modelos votan y se toma la mayoría. Para medir cuánto coinciden se usa el $\kappa$ de Fleiss, que generaliza el acuerdo corregido por azar a más de dos anotadores:
\begin{equation}
\kappa = \frac{\bar{P} - \bar{P}_e}{1 - \bar{P}_e},
\end{equation}
donde $\bar{P}$ es el acuerdo observado medio y $\bar{P}_e$ el esperado por azar. Valores cercanos a 1 indican acuerdo casi total. % TODO: \cite{TODO} Fleiss (1971) y Landis y Koch (1977) si se citan; no están en el .bib.
Un $\kappa$ alto entre modelos no implica que sus etiquetas sean correctas respecto al criterio del dominio; el Capítulo~\ref{cap:resultados} documenta un caso en el que ocurre lo contrario. Por ello se complementa con revisión de expertos y se reportan intervalos de confianza de Wilson para proporciones.

\section{Notificación, recomendación e inyección de errores}\label{sec:mt_reaccion}

La reducción de la fatiga de alertas pasa por notificar solo los eventos relevantes y con contexto suficiente para actuar. La recomendación asistida por LLM complementa la alerta con una causa probable y una acción sugerida~\cite{pedroso2025_maia}. % TODO: \cite{} para recomendación de remediación asistida por LLM (candidato en notes/03_onpremise_chaos.md: Wang et al. 2025, Preprints.org; agregar al .bib tras verificar).

La inyección de fallos y la ingeniería del caos introducen perturbaciones controladas para evaluar la resiliencia del sistema y la capacidad de los modelos de detectarlas. MAIA, por ejemplo, usa Chaos Mesh para inyectar fallos en Kubernetes y medir \textit{recall} y $F_1$ del detector por tipo de fallo, e identifica tipos con menor desempeño~\cite{pedroso2025_maia}. En este trabajo se sigue una idea análoga a escala del \textit{batch}: se inyectan eventos de error simulados en el flujo de logs y se verifica si el \textit{batch} que los contiene se clasifica como error (véase la Sección~\ref{sec:met_inyeccion}).
